\documentclass[aspectratio=169,11pt,unicode]{beamer}
\usepackage{luatexja} % 日本語を美しく処理するパッケージ
\usepackage{amsmath,amssymb,amsthm}
\usepackage{bm}

% --- テーマとカラーの設定 ---
\usetheme{Madrid}
\usecolortheme{whale}

% ナビゲーションバー（右下のアイコン）を非表示
\setbeamertemplate{navigation symbols}{}

% ブロックタイトルのフォント設定
\setbeamerfont{block title}{size=\normalsize, weight=\boldseries}

% --- タイトル情報 ---
\title[Galoisの基本定理]{Artinの定理を経由する\\Galoisの基本定理の証明}
\subtitle{講義用スライド資料}
\author{代数学講義}
\date{\today}

\begin{document}

% --- タイトルページ ---
\frame{\titlepage}

% --- 目次 ---
\begin{frame}{目次}
    \tableofcontents
\end{frame}

% =================================================================
\section{1. 基本概念の定義と例}
% =================================================================

\begin{frame}{定義 1.1 (体の拡大と自己同型群)}
\begin{block}{定義 1.1}
\begin{itemize}
    \item 体 $L$ が体 $K$ を部分体として含むとき、$L$ を $K$ の\textbf{拡大体}といい、$L/K$ と表す。
    \item $L$ は $K$ 上のベクトル空間であり、その次元 $[L:K]$ を\textbf{拡大次数}と呼ぶ。$[L:K] < \infty$ のとき\textbf{有限次拡大}と呼ぶ。
    \item 体 $L$ の自己同型全体のなす群を $\mathrm{Aut}(L)$ と書く。
    \item 部分体 $K$ の元を全て固定する自己同型のなす部分群を $\mathrm{Aut}(L/K)$ または $\mathrm{Gal}(L/K)$ と書き、$L/K$ の\textbf{Galois群}と呼ぶ。
    \item 有限群 $G \subset \mathrm{Aut}(L)$ に対して、 $L^G = \{x \in L \mid \forall \sigma \in G, \sigma(x) = x\}$ を $G$ の\textbf{不変体}と呼ぶ。
\end{itemize}
\end{block}
\end{frame}

\begin{frame}{定義 1.2 (分離性と正規性、およびGalois拡大)}
\begin{block}{定義 1.2}
体 $K$ 上の代数的な元 $\alpha$ の $K$ 上の最小多項式を考える。
\begin{itemize}
    \item \textbf{分離的 (separable):} 最小多項式が代数閉包において重根を持たないとき、$\alpha$ は $K$ 上分離的という。すべての元が分離的な拡大を\textbf{分離拡大}と呼ぶ。
    \item \textbf{正規 (normal):} $K$ 上の任意の既約多項式が $L$ に根を持つならば、その多項式が $L$ 上で一次式の積に完全に分解されるとき、$L/K$ を\textbf{正規拡大}と呼ぶ。
    \item \textbf{Galois拡大 (Galois extension):} 有限次拡大 $L/K$ が分離拡大かつ正規拡大であるとき、これを\textbf{有限次Galois拡大}と呼ぶ。
\end{itemize}
\end{block}
\end{frame}

\begin{frame}{命題 1.3 (Galois群の共役元への推移的作用)}
\begin{block}{命題 1.3}
$L/K$ が有限次Galois拡大であるとする。任意の $\alpha \in L$ と、$\alpha$ の $K$ 上での任意の共役元 $\beta \in L$（すなわち $\alpha$ の $K$ 上の最小多項式の根）に対して、ある $\sigma \in \mathrm{Gal}(L/K)$ が存在して
\[ \sigma(\alpha) = \beta \]
となる。
\end{block}
\end{frame}

\begin{frame}{命題 1.3 の証明 (1/2)}
\begin{proof}
$L/K$ は有限次分離拡大であるから、\textbf{原始元定理}により $L = K(\theta)$ となる元 $\theta \in L$ が存在する。\\
$\alpha$ と $\beta$ は $K$ 上共役であるため同じ最小多項式 $p(x) \in K[x]$ を持ち、$\alpha \mapsto \beta$ と写し $K$ の元を固定する同型写像 $\phi: K(\alpha) \xrightarrow{\sim} K(\beta)$ が一意に存在する。\\
ここで、$\theta$ の $K(\alpha)$ 上の最小多項式を $h_1(x) \in K(\alpha)[x]$ とする。$\phi$ によって $h_1(x)$ の各係数を $K(\beta)$ の元に写して得られる多項式を $h_2(x) \in K(\beta)[x]$ とおく。\\
$\theta$ の $K$ 上の最小多項式 $F(x) \in K[x]$ を考えると、$F(\theta)=0$ であるため、$h_1(x)$ は $K(\alpha)[x]$ において $F(x)$ の約数である（$F(x) = h_1(x)q_1(x)$）。$F(x)$ の係数は $K$ に属するため $\phi$ によって不変であり、したがって $h_2(x)$ も $K(\beta)[x]$ において $F(x)$ の約数となる。
\end{proof}
\end{frame}

\begin{frame}{命題 1.3 の証明 (2/2)}
\begin{proof}[証明（続き）]
$L/K$ は正規拡大であるため、$K$ 上の多項式 $F(x)$ は $L$ 内で一次式に完全に分解する。よってその約数である $h_2(x)$ も $L$ 内に根を持つ。その根の一つを $\theta' \in L$ とする。\\
多項式環の剰余環としての単拡大の構成定理より、以下の自然な同型の連鎖が得られる：
\[ L = K(\alpha)(\theta) \cong K(\alpha)[x]/(h_1(x)) \xrightarrow{\tilde{\phi}} K(\beta)[x]/(h_2(x)) \cong K(\beta)(\theta') \]
$\theta' \in L$ ゆえ $K(\beta)(\theta') \subseteq L$ であり、上の同型によって $[K(\beta)(\theta') : K] = [L : K]$ であるから、$K(\beta)(\theta') = L$ となる。\\
この連鎖を合成して得られる同型写像 $\sigma: L \xrightarrow{\sim} L$ は、$K$ の元を固定し、$\theta \mapsto \theta'$、および $\alpha \mapsto \beta$ を満たす。
したがって、$\sigma \in \mathrm{Gal}(L/K)$ であり、$\sigma(\alpha) = \beta$ が直接的に示された。
\end{proof}
\end{frame}

\begin{frame}{例 1.4 (複素数体と実数体の拡大)}
\begin{block}{例 1.4}
複素数体 $\mathbb{C}$ と実数体 $\mathbb{R}$ の拡大 $\mathbb{C}/\mathbb{R}$ を考える。
\begin{itemize}
    \item $\mathbb{C}$ の任意の元 $a+bi$ の $\mathbb{R}$ 上の最小多項式は高々2次であり、虚数部がゼロでなければ $x^2 - 2ax + a^2 + b^2$ となる。
    \item これは重根を持たず（分離的）、かつ $\mathbb{C}$ 上で完全に一次式に分解される（正規）。したがって $\mathbb{C}/\mathbb{R}$ はGalois拡大。
    \item $\mathbb{C}$ 上の自己同型で $\mathbb{R}$ を固定するものは、恒等写像 $\mathrm{id}$ と複素共役 $\sigma(a+bi) = a-bi$ のみ。
    \item したがって $\mathrm{Gal}(\mathbb{C}/\mathbb{R}) = \{\mathrm{id}, \sigma\}$ であり、不変体 $\mathbb{C}^{\{\mathrm{id}, \sigma\}}$ は $\mathbb{R}$ に一致する。
\end{itemize}
\end{block}
\end{frame}

% =================================================================
\section{2. Dedekindの補題と準備の補題}
% =================================================================

\begin{frame}{補題 2.1}
\begin{block}{補題 2.1}
体 $\Omega$ 上のベクトル空間 $V$ の $v_1, \dots, v_n \in V$ で張られる部分空間を $W$ と書く。$i \neq j$ に対して $W$ の一次変換 $T_{ij}$ が与えられており、各 $v_k \neq 0$ は、固有値 $\lambda^{ij}_k$ を持つ $T_{ij}$ の固有ベクトルであって、$\lambda^{ij}_i \neq \lambda^{ij}_j$ を満たしているとする。このとき以下が成立する。
\begin{enumerate}
    \item $i=1,\dots,n$ に対して $W$ の一次変換 $L_i$ を次で定めると、 $L_i v_j = \delta_{ij} v_i$ である。
    \[ L_i = \prod_{j \in \{1,\dots,n\} \smallsetminus \{i\}} \frac{T_{ij} - \lambda^{ij}_j \mathrm{id}_W}{\lambda^{ij}_i - \lambda^{ij}_j} \]
    \item $v_1, \dots, v_n$ は $\Omega$ 上一次独立である。
\end{enumerate}
\end{block}
\end{frame}

\begin{frame}{補題 2.1 の証明}
\begin{proof}
(1) $(T_{ij} - \lambda^{ij}_j \mathrm{id}_W) v_k = (\lambda^{ij}_k - \lambda^{ij}_j) v_k$ である。\\
$k \neq i$ ならば、積 $L_i$ の因子の中に $j=k$ となる項が存在し、分子が $(\lambda^{ik}_k - \lambda^{ik}_k) v_k = 0$ となるため、$L_i v_k = 0$ である。\\
$k = i$ の場合、すべての $j \neq i$ に対して分子は $(\lambda^{ij}_i - \lambda^{ij}_j) v_i$ となり、分母と相殺される。したがって $L_i v_i = v_i$。以上より $L_i v_j = \delta_{ij} v_i$。\\
\vspace{0.1cm}
(2) $\sum_{j=1}^n c_j v_j = 0$ ($c_j \in \Omega$) とする。この両辺に $L_i$ を作用させると、
\[ L_i \left( \sum_{j=1}^n c_j v_j \right) = \sum_{j=1}^n c_j L_i v_j = c_i v_i = 0 \]
$v_i \neq 0$ より $c_i = 0$ を得る。これが任意の $i$ で成り立つため一次独立。
\end{proof}
\end{frame}

\begin{frame}{定理 2.2 (Dedekindの補題)}
\begin{block}{定理 2.2 (Dedekindの補題)}
半群 $H$ から体 $\Omega$ への相異なる半群準同型 $\sigma_1, \dots, \sigma_n : H \to \Omega^\times$ は、$H$ 上の $\Omega$ に値を持つ関数達として一次独立である。ただし $\Omega^\times = \Omega \smallsetminus \{0\}$。
\end{block}
\end{frame}

\begin{frame}{Dedekindの補題の証明}
\begin{proof}
$V$ を写像空間、 $v_k = \sigma_k \in V$ ($v_k \neq 0$)、 $W = \langle v_1, \dots, v_n \rangle$ とする。\\
$\sigma_i$ は相異なるので、$\forall i \neq j, \exists h_{ij} \in H \text{ s.t. } \sigma_i(h_{ij}) \neq \sigma_j(h_{ij})$。\\
$W$ 上の一次変換 $T_{ij}$ を $(T_{ij} f)(x) = f(h_{ij} x)$ で定義する。\\
準同型の性質から次が成り立つ：
\[ (T_{ij} \sigma_k)(x) = \sigma_k(h_{ij} x) = \sigma_k(h_{ij}) \sigma_k(x) \]
すなわち $T_{ij} v_k = \sigma_k(h_{ij}) v_k$ であり、$v_k$ は固有値 $\lambda^{ij}_k = \sigma_k(h_{ij})$ を持つ固有ベクトル。\\
$h_{ij}$ の選び方から $\lambda^{ij}_i \neq \lambda^{ij}_j$。補題 2.1 より、 $\sigma_1, \dots, \sigma_n$ は一次独立。
\end{proof}
\end{frame}

\begin{frame}{注意 (体の準同型とDedekindの補題)}
\begin{block}{注意}
体 $L$ から体 $\Omega$ への体の準同型は、積の構造を保つため、$L$ の乗法群 $H = L^\times$ から $\Omega$ の乗法群 $\Omega^\times$ への半群準同型（群準同型）を与えます。\\
\vspace{0.3cm}
したがって、体 $L$ から体 $\Omega$ への相異なる体の準同型たちは、\textbf{Dedekindの補題}によって、直ちに関数達として $\Omega$ 上一次独立であることが従います。
\end{block}
\end{frame}

% =================================================================
\section{3. トレース写像の非退化性}
% =================================================================

\begin{frame}{命題 3.1 (トレース写像の非退化性)}
\begin{block}{命題 3.1}
体 $L$、自己同型群 $G = \{\sigma_1, \dots, \sigma_n\}$、不変体 $K = L^G$ を考える。\\
トレース写像 $\mathrm{Tr}_{L/K} : L \to K$ を $\mathrm{Tr}_{L/K}(x) = \sum_{i=1}^n \sigma_i(x)$ で定義する（これは $K$ 線形写像）。\\
\vspace{0.2cm}
このとき、すべての $y \in L$ について $\mathrm{Tr}_{L/K}(xy) = 0$ となる $x \in L$ は $x=0$ しかない。
\end{block}
\end{frame}

\begin{frame}{命題 3.1 の証明}
\begin{proof}
ある $x \in L$ について、任意の $y \in L$ に対して $\mathrm{Tr}_{L/K}(xy) = 0$ と仮定。
\[ \mathrm{Tr}_{L/K}(xy) = \sum_{i=1}^n \sigma_i(xy) = \sum_{i=1}^n \sigma_i(x)\sigma_i(y) = 0 \]
これは $L^\times$ 上の関数としての一次関係式 $\sum_{i=1}^n \sigma_i(x) \sigma_i = 0$ を意味する。\\
Dedekindの補題より、相異なる自己同型 $\sigma_1, \dots, \sigma_n$ は $L$ 上一次独立である。\\
したがって、すべての係数がゼロでなければならないため、各 $i$ に対して $\sigma_i(x) = 0$ となる。\\
$\sigma_i$ は同型写像であるから、これを満たすのは $x = 0$ のみである。
\end{proof}
\end{frame}

% =================================================================
\section{4. Artinの定理の証明}
% =================================================================

\begin{frame}{定理 4.1 (Artinの定理)}
\begin{block}{定理 4.1 (Artinの定理)}
体 $L$ とその自己同型の位数 $n$ の有限群 $G$、および不変体 $K = L^G$ を考える。このとき
\[ [L:K] = n \]
であり、$L/K$ はGalois拡大である。さらに $G = \mathrm{Gal}(L/K)$ となる。
\end{block}
\end{frame}

\begin{frame}{Artinの定理の証明 (1)}
\begin{proof}
(1) まず $\dim_L \mathrm{End}_K(L) = [L:K]$ を示す。\\
$\mathrm{End}_K(L)$ は $K$ 上の線形自己準同型環であり、左からの積作用により $L$ 上のベクトル空間とみなせる。
\begin{itemize}
    \item $[L:K] = m < \infty$ ならば、$\mathrm{End}_K(L) \cong M_m(K)$（行列環）であり、$\dim_K \mathrm{End}_K(L) = m^2$。よって $\dim_L \mathrm{End}_K(L) = m^2 / m = m = [L:K]$。
    \item $[L:K] = \infty$ ならば、$\dim_L \mathrm{End}_K(L) = \infty$ となり成立。
\end{itemize}
\end{proof}
\end{frame}

\begin{frame}{Artinの定理の証明 (2) - 単射性}
\begin{proof}[証明（続き）]
(2) 自然な写像 $\Phi: L[G] \to \mathrm{End}_K(L)$ が同型になることを示す。\\
接合環 $L[G]$ の元 $\sum_{\sigma \in G} a_\sigma \sigma$ に対し、$\Phi(\sum a_\sigma \sigma)(x) = \sum a_\sigma \sigma(x)$ とする。
\begin{itemize}
    \item $\Phi$ は環準同型である（構造から容易に確認できる）。
    \item $\Phi(\sum_{\sigma \in G} a_\sigma \sigma) = 0$ とすると、任意の $x$ で $\sum a_\sigma \sigma(x) = 0$。Dedekindの補題（一次独立性）より、すべての $\sigma$ で $a_\sigma = 0$ となり、\textbf{$\Phi$ は単射}。
\end{itemize}
\end{proof}
\end{frame}

\begin{frame}{Artinの定理の証明 (2) - 全射性 (Galois降下)}
\begin{proof}[証明（続き）]
写像 $\Psi: L \otimes_K L \to \mathrm{End}_K(L)$ を $\Psi(a \otimes b)(x) = a \mathrm{Tr}_{L/K}(bx)$ と定義。
\[ \Psi(a \otimes b) = a \sum_{\sigma \in G} \sigma(b) \sigma \in \Phi(L[G]) \implies \mathrm{Im}(\Psi) \subset \mathrm{Im}(\Phi) \]
$\sum_{i=1}^r a_i \otimes b_i \in \ker(\Psi)$ とし、$\{a_i\}$ を $K$ 上一次独立とする。\\
$\forall x, \sum a_i \mathrm{Tr}_{L/K}(b_i x) = 0 \implies \forall i, \mathrm{Tr}_{L/K}(b_i x) = 0$ ($a_i$ の一次独立性より)。\\
トレースの非退化性（命題3.1）から $b_i = 0$ となり、$\Psi$ は単射。\\
$\dim_L (L \otimes_K L) = [L:K] = \dim_L \mathrm{End}_K(L)$（有限次元の場合）より、$\Psi$ は全射（同型）。\\
したがって $\mathrm{End}_K(L) = \mathrm{Im}(\Psi) \subset \mathrm{Im}(\Phi) \subset \mathrm{End}_K(L)$ から $\Phi$ も全射。\\
次元を比較して、$[L:K] = \dim_L L[G] = |G| = n$。
\end{proof}
\end{frame}

\begin{frame}{Artinの定理の証明 (3) - Galois拡大の導出}
\begin{proof}[証明（続き）]
(3) $L/K$ が分離的かつ正規拡大であることを示す。\\
任意の $\alpha \in L$ に対し、群作用による軌道を $\{\alpha_1, \dots, \alpha_r\}$ ($\alpha_1=\alpha$) とする。\\
多項式 $f(x) = \prod_{i=1}^r (x - \alpha_i)$ は $G$ の作用で不変。係数は $L^G = K$ に属するので $f(x) \in K[x]$。\\
$\alpha$ の $K$ 上の最小多項式 $p(x)$ は $f(x)$ を割り切る。\\
$f(x)$ は重根を持たず（分離的）、$L$ 内で完全に分解する（正規）。その約数である $p(x)$ も分離的かつ正規。\\
よって $L/K$ はGalois拡大。また $[L:K] = n = |G| \le |\mathrm{Gal}(L/K)| \le [L:K]$ より $G = \mathrm{Gal}(L/K)$。
\end{proof}
\end{frame}

% =================================================================
\section{5. Galoisの基本定理の証明}
% =================================================================

\begin{frame}{定理 5.1 (Galoisの基本定理)}
\begin{block}{定理 5.1 (Galoisの基本定理)}
$L/K$ を有限次Galois拡大、 $G = \mathrm{Gal}(L/K)$ とする。\\
中間体 $M$ ($K \subset M \subset L$) と、 $G$ の部分群 $H$ の間に一対一の全単射（Galois対応）が存在する。
\[ \alpha(M) = \mathrm{Gal}(L/M), \quad \beta(H) = L^H \]
\end{block}
\end{frame}

\begin{frame}{Galoisの基本定理の証明 (1) - 包含関係}
\begin{proof}
包含関係の片側を定義から示す：
\begin{enumerate}
    \item \textbf{$M \subset \beta(\alpha(M))$ の証明：}\\
    $\alpha(M) = \mathrm{Gal}(L/M)$ は $M$ を固定する群。よって $\forall x \in M$ は $\mathrm{Gal}(L/M)$ で固定されるため、 $x \in L^{\mathrm{Gal}(L/M)} = \beta(\alpha(M))$。
    \item \textbf{$H \subset \alpha(\beta(H))$ の証明：}\\
    $\beta(H) = L^H$ は $H$ で固定される体。よって $H$ の任意の元 $\sigma$ は $L^H$ を固定するため、 $\sigma \in \mathrm{Gal}(L/L^H) = \alpha(\beta(H))$。
\end{enumerate}
\end{proof}
\end{frame}

\begin{frame}{Galoisの基本定理の証明 (2) - $H = \alpha(\beta(H))$}
\begin{proof}[証明（続き）]
\textbf{$H = \alpha(\beta(H))$ の等号：}\\
$H \subset G$ を任意の部分群とし、不変体 $M' = \beta(H) = L^H$ を考える。\\
\vspace{0.2cm}
Artinの定理（定理4.1）を体 $L$ と部分群 $H$ に対し適用する。
\begin{itemize}
    \item $[L : L^H] = |H|$ となる。
    \item さらに $L/L^H$ はGalois拡大であり、そのGalois群は $H$ に一致する。
\end{itemize}
\vspace{0.2cm}
すなわち、 $\mathrm{Gal}(L/L^H) = H$。\\
これは定義より $\alpha(\beta(H)) = H$ を意味する。
\end{proof}
\end{frame}

\begin{frame}{Galoisの基本定理の証明 (3) - $M = \beta(\alpha(M))$}
\begin{proof}[証明（続き）]
\textbf{$M = \beta(\alpha(M))$ の等号：}\\
$M' = \beta(\alpha(M))$ とおく。 $M \subset M'$ は既知。\\
先ほどの結果に $H = \alpha(M)$ を代入すると、 $\alpha(M') = \alpha(M)$ を得る。\\
もし $M \subsetneq M'$ ならば、ある $\gamma \in M' \smallsetminus M$ が存在。\\
$L/M$ も有限次Galois拡大であり、 $\gamma$ の $M$ 上の最小多項式 $p(x)$ は $\deg p \ge 2$ かつ分離的。よって $\gamma \neq \gamma'$ なる共役元 $\gamma' \in L$ を持つ。\\
ここで\textbf{命題 1.3} を $L/M$ に適用すると、 $\sigma(\gamma) = \gamma'$ となる $\sigma \in \mathrm{Gal}(L/M) = \alpha(M)$ が存在する。\\
しかし $\alpha(M) = \alpha(M')$ なので、 $\sigma$ は $M'$ の元 $\gamma$ を固定しなければならず、 $\sigma(\gamma) = \gamma$。これは矛盾。\\
したがって $M = M' = \beta(\alpha(M))$。
\end{proof}
\end{frame}

\begin{frame}{参考文献}
\begin{itemize}
    \item Artin, E. (1971). \textit{Galois Theory: Lectures Delivered at the University of Notre Dame} (Notre Dame Mathematical Lectures, Number 2). Project Euclid.
    \item Lang, S. (2002). \textit{Algebra} (Revised 3rd ed.). Springer.
\end{itemize}
\end{frame}

\end{document}